\subsection{مبانی نظری و فرمول‌بندی ریاضی}
در مرحله کالیبراسیون، ماتریس کوواریانس تجمعی فعال‌سازی‌های ورودی $\mathbf{C}_X \in \mathbb{R}^{d_{\text{in}} \times d_{\text{in}}}$ برای هر لایه استخراج و با مقادیر ویژه منظم‌سازی می‌گردد:
\begin{equation}
\mathbf{C}_X = \mathbf{Q} \mathbf{\Lambda} \mathbf{Q}^\top, \quad \mathbf{C}_X^{1/2} = \mathbf{Q} \mathbf{\Lambda}^{1/2} \mathbf{Q}^\top, \quad \mathbf{C}_X^{-1/2} = \mathbf{Q} \mathbf{\Lambda}^{-1/2} \mathbf{Q}^\top
\end{equation}
ماتریس وزن $\mathbf{W} \in \mathbb{R}^{d_{\text{out}} \times d_{\text{in}}}$ با جذر کوواریانس پیش‌شرط‌سازی شده و تجزیه مقادیر منفرد بر روی ماتریس سفیدشده انجام می‌پذیرد:
\begin{equation}
\mathbf{M} = \mathbf{W} \mathbf{C}_X^{1/2} = \mathbf{U} \mathbf{\Sigma} \mathbf{V}^\top
\end{equation}
طیف مقادیر منفرد $\mathbf{\Sigma} = \operatorname{diag}(\sigma_1, \dots, \sigma_d)$ با نرمال‌سازی به یک توزیع احتمالاتی بر روی زیرفضاها تبدیل می‌شود:
\begin{equation}
p_i = \frac{\sigma_i}{\sum_{j=1}^d \sigma_j + \epsilon}
\end{equation}
آنتروپی شانون طیف در مقیاس بیت به فرم زیر محاسبه می‌گردد:
\begin{equation}
H(\mathbf{\Sigma}) = -\sum_{i=1}^d p_i \log_2(p_i)
\end{equation}
برای هر مؤلفه کاندید $i \in [r_{\text{floor}}, r_{\max}]$، امتیاز اهمیت اطلاعاتی $\mathcal{S}_{l, i}$ تعریف می‌شود:
\begin{equation}
\mathcal{S}_{l, i} = \frac{\sigma_{l, i}^2}{H(\mathbf{\Sigma}_l)}
\end{equation}
تمام مؤلفه‌های طیفی در سراسر لایه‌ها با امتیاز $\mathcal{S}_{l, i}$ مرتب شده و باقیمانده بودجه سراسری به برترین مؤلفه‌ها تخصیص می‌یابد تا رتبه بهینه $r_l$ تعیین گردد. در نهایت فاکتورهای کم‌رتبه به صورت متقارن و متعادل با معکوس جذر کوواریانس تشکیل می‌شوند:
\begin{equation}
\mathbf{U}_{r_l} = \mathbf{U}_l[:, :r_l] \mathbf{\Sigma}_l[:r_l]^{1/2}, \quad \mathbf{V}_{r_l} = (\mathbf{\Sigma}_l[:r_l]^{1/2} \mathbf{V}_l[:, :r_l]^\top) \mathbf{C}_{X, l}^{-1/2}
\end{equation}
